\section{Conclusion}

We presented OSTRA, a framework for multi-stage manipulation that represents
historical and future object states in a shared Object-State Transition Map.
Transition-Guided Attention and Staged Action Generation progressively translate
object-state changes into TCP subgoals and executable actions. Across RLBench2,
ManiSkill, LIBERO, and seven real-world tasks, OSTRA consistently outperforms
representative imitation-learning baselines, while the ablations validate the
contributions of OST-Map and the proposed architectural designs. Together,
these results show that explicit transition reasoning provides a structured
bridge between perception and continuous control. Future work will improve map
construction under perception errors and extend object-state reasoning to less
structured, open-world manipulation with stronger transfer across embodiments
and task domains, including deformable objects, dynamic scenes, and broader
task families.
